La risposta di un mezzo poroso saturo a una sollecitazione meccanica dipende dall'interazione fra deformazione dello scheletro, pressione interstiziale e moto del fluido. Nella formulazione di Biot~\citep{biot1941general}, la deformazione modifica il contenuto di fluido e la pressione contribuisce allo stato di sforzo; il gradiente di pressione determina quindi un flusso che può trasportare specie disciolte. In un materiale con risposta ereditaria, tuttavia, la stessa storia di carico può produrre un'evoluzione temporale non descrivibile da un unico tempo di rilassamento. Anche il trasporto di soluti in mezzi microstrutturalmente complessi può mostrare code temporali e dinamiche non fickiane~\citep{metzler2000the,sykov2008diffusion}. Una questione centrale per la modellazione accoppiata è allora distinguere \emph{dove} sia introdotta la memoria: nella legge costitutiva dello scheletro, nella relazione di flusso o nell'equazione che governa il tracciante. Queste scelte non sono intercambiabili e conducono a differenti meccanismi di trasferimento dell'informazione verso le grandezze osservabili.

Il tessuto cerebrale fornisce un contesto motivante, pur senza costituire in questo lavoro un sistema clinicamente calibrato. Le sue proprietà meccaniche variano con regione, protocollo e scala temporale~\citep{budday2019fifty}; formulazioni bifasiche poro-viscoelastiche descrivono l'interazione fra fase solida, fluido e rilassamento viscoso~\citep{mehrabian2015dual,comellas2020modeling}. Greiner e collaboratori hanno mostrato, su dati di prove meccaniche, che i parametri porosi e viscoelastici possono compensarsi e che la loro calibrazione beneficia di osservabili meccanici complementari~\citep{greiner2024modeldriven}. Sul versante del trasporto, la diffusione e l'avvezione interstiziali restano oggetto di discussione e dipendono dal regime sperimentale~\citep{ray2019analysis,alcaide2023solute}; misure e modelli basati su traccianti dinamici sono stati usati per inferire diffusività apparenti e caratteristiche del trasporto cerebrale~\citep{valnes2020apparent}. Sviluppi recenti ricostruiscono campi di velocità da immagini dinamiche e simulano il trasferimento di traccianti in anatomie affette da glioma~\citep{toscano2026mraiv,poulain2026glioma}; obiettivi e osservabili di questi studi differiscono dalla calibrazione dei due ordini di memoria qui affrontata. Ne consegue che il solo segnale di concentrazione costituisce una misura indiretta della meccanica: dipende tanto dalle proprietà del trasporto quanto dal flusso prodotto dalla deformazione.

\subsection{Stato dell'arte e collocazione della formulazione}
\label{sec:intro-literature}

Le tre principali linee di ricerca pertinenti hanno obiettivi e sedi della memoria distinti. La prima riguarda la \emph{poro-viscoelasticità}: l'estensione delle equazioni di Biot con dissipazione dello scheletro è consolidata sia nei modelli computazionali di tessuti~\citep{comellas2020modeling,greiner2024modeldriven} sia nell'analisi delle equazioni quasi-statiche con viscoelasticità classica~\citep{bociu2023mathematical}. Leggi di Kelvin--Voigt di ordine frazionario e relative proprietà dissipative sono già note in viscoelasticità~\citep{libertiaux2010differential,caponi2021sapio}; tali risultati forniscono il fondamento costitutivo e matematico della memoria dello scheletro qui adottata.

La seconda linea sviluppa \emph{modelli poromeccanici con operatori frazionari}. Alaimo e collaboratori introducono una memoria temporale nella relazione fra flusso e gradiente di pressione~\citep{alaimo2019a}; Barrera formula e discretizza mediante elementi finiti una diffusione anomala della pressione accoppiata alla meccanica~\citep{barrera2021unified}. Amiri e collaboratori descrivono invece l'ereditarietà della risposta di un tessuto biologico saturo mediante una legge costitutiva frazionaria, mantenendo il flusso di Darcy~\citep{amiri2023fractionalorder}. Tyrylgin e collaboratori considerano poroelasticità temporale frazionaria in mezzi eterogenei e fratturati~\citep{tyrylgin2023computational}. Il confronto è essenziale perché rende esplicita la differenza tra memoria idraulica e memoria dello scheletro, che nel presente modello occupano posizioni costitutive distinte.

La terza linea studia il \emph{trasporto anomalo in mezzi porosi deformabili}. La modellazione frazionaria della diffusione è ampiamente documentata~\citep{metzler2000the,magin2008anomalous}; in particolare, Salomoni e De Marchi hanno già accoppiato un modello idromeccanico per un mezzo saturo deformabile a equazioni frazionarie di flusso e trasporto di soluti~\citep{salomoni2022fractional}. La nostra analisi si concentra sull'effetto congiunto di memoria costitutiva e accumulo frazionario nella sola osservazione di concentrazione, mantenendo classica la legge di Darcy.

\begin{table*}[!t]
\centering\footnotesize
\caption{Posizionamento rispetto ai lavori più vicini. La distinzione riguarda la \emph{sede} degli operatori di memoria e gli osservabili indagati, anziché una graduatoria dei modelli. Le descrizioni sono circoscritte agli aspetti pertinenti al presente lavoro.}
\label{tab:literature-positioning}
\setlength{\tabcolsep}{4.0pt}
\renewcommand{\arraystretch}{1.13}
\begin{tabularx}{\textwidth}{@{}>{\raggedright\arraybackslash}p{0.23\textwidth}>{\raggedright\arraybackslash}p{0.30\textwidth}>{\raggedright\arraybackslash}X@{}}
\toprule
Riferimento & Sede della memoria e formulazione & Contributo pertinente e differenza rispetto a questo studio\\
\midrule
Biot~\citep{biot1941general}; Bociu et al.~\citep{bociu2023mathematical} & Consolidazione classica; viscoelasticità lineare nel sottosistema pressione--spostamento. & Fondamento dell'accoppiamento idromeccanico e dell'analisi dissipativa, senza un ordine frazionario osservato tramite tracciante.\\
Mehrabian et al.~\citep{mehrabian2015dual}; Comellas et al.~\citep{comellas2020modeling} & Scheletro viscoelastico e flusso interstiziale in modelli di tessuto cerebrale. & Motivano l'interazione fra meccanica e fluido; non costituiscono il confronto a due ordini frazionari indipendenti qui esaminato.\\
Alaimo et al.~\citep{alaimo2019a}; Barrera~\citep{barrera2021unified} & Memoria nella legge idraulica o nella diffusione di pressione. & Poromeccanica frazionaria già consolidata, ma con operatori posti in sedi differenti dalla memoria del trasporto qui scelta.\\
Amiri et al.~\citep{amiri2023fractionalorder}; Caponi e Sapio~\citep{caponi2021sapio} & Memoria costitutiva frazionaria del solido. & Base per descrivere l'ereditarietà meccanica; non affrontano lo stesso canale osservativo tracciante--flusso--meccanica.\\
Salomoni e De Marchi~\citep{salomoni2022fractional} & Flusso non darciano e trasporto frazionario di soluti in mezzo deformabile. & Precedente diretto dell'accoppiamento trasporto--deformazione, con operatori e obiettivi modellistici differenti.\\
Greiner et al.~\citep{greiner2024modeldriven}; Valnes et al.~\citep{valnes2020apparent} & Stima di parametri meccanici da prove di carico; trasporto di tracciante analizzato mediante dati dinamici. & Chiariscono, in contesti sperimentali distinti, le difficoltà della calibrazione e il carattere indiretto delle osservazioni.\\
Presente lavoro & Kelvin--Voigt--Caputo nello scheletro ($\beta$), Darcy classico, Caputo nell'accumulo del tracciante ($\alpha$); accoppiamento unidirezionale verso $c$. & Confronto fattoriale M1--M4, verifica della risposta osservabile e analisi della riconoscibilità della memoria meccanica usando sole concentrazioni sintetiche.\\
\bottomrule
\end{tabularx}
\end{table*}

\subsection{Questione aperta, obiettivi e contributi}
\label{sec:intro-gap}

Il quadro bibliografico mostra dunque che gli ingredienti principali sono già disponibili, ma vengono spesso interrogati con \emph{osservabili differenti}: il rilassamento e la pressione per la poromeccanica, la concentrazione per il trasporto, le curve di carico per la calibrazione dei parametri viscoelastici. Il problema che affrontiamo consiste nel chiarire le conseguenze della \emph{collocazione} degli operatori di memoria in un sistema accoppiato e stabilire quanto di tali conseguenze sopravviva quando si osserva soltanto il tracciante. In particolare, una variazione di $\beta$ modifica la risposta meccanica, il flusso di Darcy e infine la concentrazione; una variazione di $\alpha$ agisce direttamente sull'accumulo del tracciante. Questi due cammini possono produrre firme simili e risentire anche di $E$, $D$, $K_h$ e dell'ampiezza viscoelastica $A_v$. La distinzione fra grandezza dell'effetto sul forward e informazione parametrica indipendente è quindi parte della verifica del modello, non una proprietà desumibile dal solo rango del Jacobiano.

Formuliamo pertanto la domanda guida nel modo seguente: \emph{in un modello poro-viscoelastico con memoria frazionaria costitutiva e memoria frazionaria del tracciante, quali effetti sono imputabili ai due meccanismi, come essi interagiscono nella concentrazione e in quali condizioni la memoria meccanica rimane distinguibile dagli altri coefficienti e dalla discrepanza numerica}. L'esame ha carattere modellistico e computazionale: una sezione BraTS fornisce il dominio e le regioni, mentre le forzanti, i coefficienti materiali e le osservazioni temporali sono prescritti e sintetici~\citep{menze2015the,bakas2017advancing,brats2020protocol}.

Per rispondere a questa domanda il lavoro propone tre livelli di contributo, con la modellazione come punto di partenza. \emph{Primo}, definiamo e analizziamo una formulazione triangolare Biot--Kelvin--Voigt--Caputo--Darcy--trasporto nella quale gli ordini $\beta$ e $\alpha$ agiscono in sedi fisicamente separate. Il confronto controllato fra le quattro varianti M1--M4 isola gli effetti individuali delle memorie e permette di quantificarne la non additività osservativa. Il confronto è interpretato come analisi strutturale e non come selezione clinica del modello. \emph{Secondo}, esaminiamo il canale di osservazione meccanica--flusso--tracciante: le sensibilità locali e i profili riottimizzati distinguono il rango globale dall'informazione residua sulla memoria meccanica. L'analisi si fonda su sensibilità condizionate e profili di costo riottimizzati, strumenti consolidati nel disegno ottimale e nei problemi inversi frazionari~\citep{pukelsheim2006optimal,raue2009structural,jin2015tutorial}. \emph{Terzo}, verifichiamo il solutore mediante test manufatti e bilanci, e confrontiamo le osservazioni su discretizzazioni differenti: nei problemi inversi la discrepanza del modello diretto non può essere confusa a priori con il rumore di misura~\citep{kaipio2007discretization}. Questo terzo livello specifica la scala entro cui interpretare le differenze fra varianti e gli indizi di identificabilità: le norme delle discrepanze osservazionali misurano un effetto numerico, mentre la loro proiezione sulle sensibilità parametriche è un problema distinto.

La Sezione~\ref{sec:forward} presenta la formulazione e le sue proprietà strutturali; la Sezione~\ref{sec:inverse} definisce osservazioni e problemi di identificabilità. La Sezione~\ref{sec:numerics} descrive verifiche, varianti e protocolli, mentre la Sezione~\ref{sec:results} discute prima la risposta dei modelli e poi le sue implicazioni per l'inferenza dei parametri.
